\label{sec:results}

\subsection{National Count of Competitive Districts}

Table~\ref{tab:competitive-national} presents the key finding.  Algorithmic
maps produce 85 competitive districts (margin under 10 points in the two-party
presidential vote) compared with 65 under enacted plans---a 31\% increase.
Among the competitive districts, 57 are swing districts (margin under 5
points) under algorithmic maps versus 38 under enacted maps.

\begin{table}[h]
\centering
\caption{Competitive Districts: Algorithmic vs.\ Enacted Maps, 2020}
\label{tab:competitive-national}
\begin{threeparttable}
\begin{tabular}{lcc}
\toprule
Measure & Algorithmic & Enacted \\
\midrule
Competitive districts ($<$10 pt margin) & 85 & 65 \\
Swing districts ($<$5 pt margin) & 57 & 38 \\
Safe Democratic ($>$10 pt D margin) & 189 & 208 \\
Safe Republican ($>$10 pt R margin) & 161 & 162 \\
\midrule
Total districts & 435 & 435 \\
\midrule
Competitive share & 19.5\% & 14.9\% \\
Swing share & 13.1\% & 8.7\% \\
\bottomrule
\end{tabular}
\begin{tablenotes}
\small
\item Note: Margins computed from 2020 presidential two-party vote share
  reaggregated to district boundaries.  Safe districts defined as those with
  a two-party presidential margin exceeding 10 percentage points.
\end{tablenotes}
\end{threeparttable}
\end{table}

The increase in competitive districts (from 65 to 85) comes primarily from
a reduction in safe Democratic seats, which fall from 208 to 189.  Safe
Republican seats are nearly unchanged (162 to 161).  This asymmetry reflects
the geographic efficiency advantage enjoyed by Republicans in a single-member
district system: Democratic voters are concentrated in urban cores, producing
safe Democratic seats regardless of the map-drawing method.  Algorithmic maps
partially offset this by drawing more geographically diverse districts that
mix urban and suburban tracts, but the geographic sorting effect identified
by \citet{chen2013unintentional} and \citet{rodden2019why} is not eliminated.

\subsection{Competitive Districts Are Genuinely Contested}

The 85 competitive algorithmic districts are not artifacts of the vote-share
definition.  Examined individually:
\begin{itemize}
  \item 72 of the 85 (85\%) have presidential margins under 8 points.
  \item 57 of the 85 (67\%) have margins under 5 points.
  \item The median margin among competitive algorithmic districts is 4.1
    percentage points.
\end{itemize}

By comparison, among the 65 enacted competitive districts:
\begin{itemize}
  \item 49 of the 65 (75\%) have margins under 8 points.
  \item 38 of the 65 (58\%) have margins under 5 points.
  \item The median margin among competitive enacted districts is 5.3
    percentage points.
\end{itemize}

Algorithmic competitive districts are on average \emph{more} competitive
within the competitive window, not merely more numerous.

\subsection{The Byproduct Character of Competitiveness}

A critical qualification is in order.  The algorithm does not target
competitiveness; it targets population equality and compactness.  The
increase in competitive districts is a byproduct of geometric neutrality,
not a design feature.  Indeed, competitiveness is a criterion that a neutral
algorithm \emph{cannot} target without accessing partisan data---the very
information the algorithm is designed to ignore.

This point has legal significance.  An algorithm that targeted competitiveness
would necessarily access voter registration or election data, raising the same
concerns about partisan manipulation (inverted) that motivated algorithmic
redistricting in the first place.  The appropriate characterization of the
finding is: \emph{compactness optimization produces more competitive districts
than partisan optimization does, because partisan optimizers eliminate
competitive races by design}.
